%!TEX root = ../_fyp.tex
\chapter{Discussion}
\label{chap:discussion}

Chapter~\ref{c-results} established that Ask, Code, and Agentic modes satisfy distinct contracts on shared notebook snapshots. This chapter interprets those findings relative to prior work on notebook assistants \parencite{mcnutt2023notebookassistants,rule2018exploration}, tool-augmented \acp{LLM} \parencite{schick2023toolformer,qin2024toolllm}, and interactive programming cost models \parencite{mozannar2024reading}.

\section{Agentic Tool Dispatch}

The primary Agentic result---37 tools dispatched in the second \ac{API} call after a read-only query round---confirms that GLM~4.7 can emit large parallel \texttt{tool\_calls} batches when native tool schemas are attached \parencite{qin2024toolllm}. The two-call design trades the flexibility of open-ended ReAct trajectories \parencite{yao2023react} for predictable latency and cost, consistent with fire-and-forget batching described in Section~\ref{sec:tool-registry}.

Importantly, the present artefact \emph{does not} close the loop with per-cell observation and replanning. Dispatch success therefore indicates planning fidelity at the host boundary, not guaranteed kernel outcomes in the live editor. This distinction aligns with \textcite{hou2024llm4se}, who note that software-engineering agents require explicit evaluation criteria matched to deployment contracts. For the FYP, dispatch-only scoring is the appropriate metric until a future ReAct implementation is added (Section~\ref{sec:future-react}).

Host-side descending reorder for bulk inserts and deletes mitigated index-shift risk identified in notebook integration literature \parencite{rule2019tenrules,chasins2018roussillon}. The Pakistan housing pipeline batch suggests that domain-specific tool narrowing---six vetted cell operations rather than open web APIs \parencite{patil2023gorilla}---improves batch coherence.

\section{Ask Mode: Explanation Without Side Effects}

Ask mode achieved strong performance on cell-local explanation (cell~17) and error diagnosis (cell~31) while maintaining zero side effects. Zero tool calls eliminate the verification burden that \textcite{vaithilingam2022expectation} associate with generative coding tools, because the assistant cannot mutate the artefact under discussion.

The notebook workflow test exhibited 89\% topic coverage, indicating that the model recognised standard data-science stages (loading, cleaning, feature work, modelling, evaluation) from packed context---a form of lightweight retrieval grounding \parencite{lewis2020rag}. The remaining limitation is depth: without full notebook coverage in the prompt window \parencite{wei2022chain,brown2020language}, workflow answers may summarise stages rather than exhaustively trace every cell. This is an expected trade-off for closed-platform assistants that lack platform \ac{API} access \parencite{chasins2018roussillon}.

For educational use in \acf{BSDS} coursework, Ask mode is best positioned as a \emph{read-only tutor}: explain cells, diagnose errors, and clarify intent before the student authors or runs code \parencite{finnieansley2022codex}.

\section{Code Mode: Implementation Guidance}

Code mode passed all four reported scenarios while preserving the no-write contract. Correct XGBoost training code (100\% keyword alignment with notebook variables) shows that snapshot grounding suffices for variable reuse without live kernel introspection \parencite{lewis2020rag}. Cell replacement and feature-engineering tests further indicate that the model can emit self-contained cells matching notebook naming conventions \parencite{wang2021welldocumented}.

The empty-cell workflow result---asking intent before emitting Placement and Code sections---implements the deferral pattern recommended for notebook copilots \parencite{mcnutt2023notebookassistants} and reduces the risk of unsolicited code dumps \parencite{mozannar2024reading}.

The principal weakness is placement accuracy under partial context: correct code may be paired with suboptimal insert indices. This mirrors IDE copilot findings that suggestion quality depends on visible context \parencite{ziegler2024copilot}, extended to positional notebook cells rather than flat files.

\section{Mode Selection Implications}

Table~\ref{tab:mode-guidance} translates results into practical guidance. The separation supports Objective~6 in Section~\ref{sec:objectives}: empirical validation across modes, not a single monolithic assistant score.

\begin{table}[!ht]
\centering
\caption{Practical mode selection guidance derived from benchmarks.}
\label{tab:mode-guidance}
\small
\begin{tabular}{|p{3.2cm}|p{4.2cm}|p{4.8cm}|}
\hline
\textbf{User goal} & \textbf{Recommended mode} & \textbf{Evidence} \\ \hline
Understand notebook or cell & Ask & 89\% workflow coverage; 0\% hallucination on cell~17 \\
Diagnose execution error & Ask & Correct root-cause on cell~31 \\
Implement new cell safely & Code & Runnable blocks; 0 browser writes \\
Automate multi-cell batch & Agentic & 37-tool dispatch after query call \\ \hline
\end{tabular}
\end{table}

\section{Threats to Validity}

Three threats qualify the generalisability of findings. \textbf{Harness fidelity:} Agentic tests mock browser execution; traces record dispatch, not live Kaggle latency. \textbf{Heuristic scoring:} Coverage and code-correctness metrics automate keyword rubrics rather than expert judgment \parencite{chen2021codex}. \textbf{Single-model scope:} All runs use GLM~4.7; cross-model replication would strengthen claims \parencite{hou2024llm4se}. Despite these limits, the benchmarks provide reproducible evidence that mode contracts are enforceable in a deployable Chrome extension architecture.

\section{Comparison with Prior Notebook Assistants}
\label{sec:comparison-prior}

\textcite{mcnutt2023notebookassistants} prototyped inline suggestions within JupyterLab-style environments where extension APIs expose cell metadata directly. The present artefact operates on Kaggle's closed \texttt{/edit} surface without platform cooperation, trading API richness for \ac{DOM} scraping \parencite{chasins2018roussillon}. Relative to their design probes, this FYP adds three distinctions:

\begin{itemize}
  \item \textbf{Explicit mode contracts} separate read-only tutoring (Ask), copy-paste generation (Code), and batched dispatch (Agentic) rather than a single inline suggestion channel.
  \item \textbf{Native tool calling} replaces free-form code fences with schema-validated \texttt{insert\_cell} and \texttt{run\_cell} operations \parencite{qin2024toolllm}.
  \item \textbf{Host-side batch reordering} addresses positional index shift---a notebook-specific hazard absent in flat-file copilots \parencite{rule2019tenrules}.
\end{itemize}

IDE copilots such as GitHub Copilot excel at file-buffer completion \parencite{ziegler2024copilot} but do not model markdown/code interleaving or kernel execution order. The benchmark evidence suggests that notebook assistance requires platform-specific grounding even when the underlying model is shared.

\section{Educational Implications for BSDS}
\label{sec:edu-implications}

For \acf{BSDS} coursework at \acf{GCUF}, the mode separation maps onto pedagogical stages:

\begin{enumerate}
  \item \textbf{Exploration phase.} Students use Ask mode to understand forked Kaggle kernels before editing---mirroring \textcite{rule2018exploration}'s distinction between exploration and explanation notebooks.
  \item \textbf{Implementation phase.} Code mode supports assessed assignments where students must author cells themselves but benefit from placement hints and runnable templates \parencite{finnieansley2022codex}.
  \item \textbf{Automation phase.} Agentic mode demonstrates how tool-augmented models plan multi-cell workflows, preparing advanced students for agentic data-science tooling while remaining bounded by the two-call contract.
\end{enumerate}

The empty-cell deferral results (Code Test~4) are particularly relevant: assistants should clarify intent before filling cells in graded work, reducing academic integrity concerns associated with unsolicited code dumps \parencite{vaithilingam2022expectation}.

\section{Architectural Trade-offs}
\label{sec:arch-tradeoffs}

The split extension--host design embodies deliberate trade-offs summarised in Table~\ref{tab:tradeoffs}.

\begin{table}[!ht]
\centering
\caption{Architectural trade-offs in the implemented system.}
\label{tab:tradeoffs}
\small
\begin{tabular}{|p{3.4cm}|p{4.8cm}|p{4.8cm}|}
\hline
\textbf{Decision} & \textbf{Benefit} & \textbf{Cost} \\ \hline
Native messaging & API keys never in extension; full-duplex I/O & Local install and host registration \\ \hline
DOM scraping & Works without Kaggle API & Fragile to front-end changes \\ \hline
Two-call Agentic cap & Predictable latency and billing & No mid-batch replanning \\ \hline
Fire-and-forget dispatch & High throughput (37 tools in one batch) & No execution verification in same turn \\ \hline
Dual JSON stores & Fresh live + durable persistent context & Storage sync complexity \\ \hline
\end{tabular}
\end{table}

These trade-offs reflect FYP scope constraints (Section~\ref{sec:scope}) while remaining transferable to other closed notebook platforms identified in Chapter~\ref{c-review}.

\section{Context Packing as a Cross-Mode Bottleneck}
\label{sec:context-bottleneck}

Ask Test~1 and Code placement limitations share a root cause: the host packs notebook context subject to \texttt{MAX\_NOTEBOOK\_CONTEXT\_CHARS} and \texttt{MAX\_INPUT\_TOKENS}. Long competition notebooks exceed single-window coverage \parencite{wei2022chain}. The model's appropriate response---high-level workflow summary rather than per-cell invention---is pedagogically sound but reduces automated accuracy scores.

Mitigations proposed in Section~\ref{sec:future-react} include hierarchical section summaries and retrieval over cell embeddings \parencite{lewis2020rag,jiang2023structgpt}. Until implemented, practitioners should prefer cell-local Ask queries and verify Code placement against the visible editor.

\section{Reliability and Safety Posture}
\label{sec:safety}

From a safety perspective, mode contracts provide defence in depth:

\begin{itemize}
  \item Ask mode cannot mutate notebooks---eliminating accidental destructive edits during explanation.
  \item Code mode requires explicit user copy-paste---preserving human approval before execution.
  \item Agentic mode batches all writes in one host-validated queue---enabling user review of planned operations before browser execution completes.
\end{itemize}

This posture aligns with \textcite{mozannar2024reading}'s finding that users accept AI assistance when edit semantics are predictable. Future ReAct loops must preserve these guardrails rather than bypassing them for autonomy \parencite{shinn2023reflexion}.
